% Automatically generated by scripts/generate_report_figures.py.
% Do not edit manually.

\newcommand{\HuFinalTemperature}{856.8}
\newcommand{\HuIntegratedEnergy}{386.6}
\newcommand{\HuSeiEnergy}{9.4}
\newcommand{\HuSeiFraction}{2.4}
\newcommand{\HuSeiFinalConversion}{1.000}
\newcommand{\HuAnodeEnergy}{314.0}
\newcommand{\HuAnodeFraction}{81.2}
\newcommand{\HuAnodeFinalConversion}{1.000}
\newcommand{\HuCathodeEnergy}{37.9}
\newcommand{\HuCathodeFraction}{9.8}
\newcommand{\HuCathodeFinalConversion}{1.000}
\newcommand{\HuElectrolyteEnergy}{25.2}
\newcommand{\HuElectrolyteFraction}{6.5}
\newcommand{\HuElectrolyteFinalConversion}{1.000}
\newcommand{\HuSensibleEnergy}{376.8}
\newcommand{\HuConvectiveLoss}{9.6}
\newcommand{\HuEnergyResidual}{0.1563}
\newcommand{\HuRelativeEnergyResidual}{0.0404}


The behavior of the implemented thermal runaway framework is illustrated
using a reference simulation based on the reaction network of
Hu et al.~\cite{hu2020inhibition}.

The purpose of this simulation is to demonstrate the coupled response of
the implemented reaction and thermal models. The calculation is not
intended as an independent experimental validation of the original model.


\subsection{Reference simulation conditions}

The simulation is initialized at an elevated temperature in order to
activate the thermal runaway reaction network within a short numerical
time interval.

The principal simulation conditions are summarized in
\cref{tab:reference_simulation_conditions}.

\begin{table}[htbp]
    \centering
    \caption{Conditions used for the Hu et al.\ reference simulation.}
    \label{tab:reference_simulation_conditions}

    \begin{tabular}{lll}
        \toprule
        Parameter & Value & Unit \\
        \midrule
        Initial temperature
        & 480
        & \si{\kelvin}
        \\

        Simulation duration
        & 20
        & \si{\second}
        \\

        Convective heat-transfer coefficient
        & 10
        & \si{\watt\per\meter\squared\per\kelvin}
        \\

        Maximum permitted temperature
        & 1200
        & \si{\kelvin}
        \\

        \bottomrule
    \end{tabular}
\end{table}


\subsection{Temperature evolution}

The simulated temperature evolution is shown in
\cref{fig:hu_temperature_result}.

\begin{figure}[htbp]
    \centering
    \includegraphics[
        width=0.78\textwidth
    ]{hu2020_temperature.pdf}

    \caption{
        Cell temperature predicted by LiionTR for the Hu et al.\
        reference thermal runaway simulation.
    }
    \label{fig:hu_temperature_result}
\end{figure}


\subsection{Reaction progression}

The evolution of the four reaction conversions is shown in
\cref{fig:hu_conversion_result}.

\begin{figure}[htbp]
    \centering
    \includegraphics[
        width=0.82\textwidth
    ]{hu2020_conversion.pdf}

    \caption{
        Evolution of the SEI, anode-electrolyte, cathode, and electrolyte
        reaction conversions during the reference simulation.
    }
    \label{fig:hu_conversion_result}
\end{figure}


\subsection{Reaction heat generation}

The individual reaction contributions to heat generation are shown in
\cref{fig:hu_heat_result}.

\begin{figure}[htbp]
    \centering
    \includegraphics[
        width=0.82\textwidth
    ]{hu2020_heat_generation.pdf}

    \caption{
        Mass-specific heat-generation rates associated with the thermal
        runaway reactions in the Hu et al.\ reference simulation.
    }
    \label{fig:hu_heat_result}
\end{figure}

\subsection{Reaction-wise energy release}

Integration of the individual heat-generation rates provides the
specific energy released by each reaction during the reference
simulation.

The resulting energy balance is summarized in
\cref{tab:hu_reaction_summary}.

\input{generated/hu2020_reaction_summary}

\subsection{Energy distribution}

The energy distribution in
\cref{tab:hu_reaction_summary}
shows that the anode--electrolyte reaction is the dominant energetic
contribution in the present reference simulation.

It releases approximately

\begin{equation}
    q_{\mathrm{an}}
    =
    \SI{\HuAnodeEnergy}{\kilo\joule\per\kilogram},
\end{equation}

corresponding to approximately

\begin{equation}
    \HuAnodeFraction\%
\end{equation}

of the total reaction energy.

The cathode decomposition represents the second-largest contribution,

\begin{equation}
    q_{\mathrm{cat}}
    =
    \SI{\HuCathodeEnergy}{\kilo\joule\per\kilogram},
\end{equation}

or approximately

\begin{equation}
    \HuCathodeFraction\%
\end{equation}

of the total.

The electrolyte decomposition contributes

\begin{equation}
    q_{\mathrm{el}}
    =
    \SI{\HuElectrolyteEnergy}{\kilo\joule\per\kilogram},
\end{equation}

while SEI decomposition contributes only

\begin{equation}
    q_{\mathrm{SEI}}
    =
    \SI{\HuSeiEnergy}{\kilo\joule\per\kilogram}.
\end{equation}

The relatively small integrated SEI contribution is partly explained
by its initial reaction state. The SEI conversion is initialized at

\begin{equation}
    \alpha_{\mathrm{SEI},0}
    =
    0.85,
\end{equation}

so only a relatively small remaining fraction is available for further
decomposition during the simulated interval.

In contrast, the anode--electrolyte reaction begins at

\begin{equation}
    \alpha_{\mathrm{an},0}
    =
    0.25,
\end{equation}

leaving a substantially larger available reaction extent together with
a high specific reaction enthalpy.


\subsection{Reaction completion}

All four reaction states approach complete conversion during the
reference simulation,

\begin{equation}
    \alpha_i
    \rightarrow
    1.
\end{equation}

Consequently, the differences in integrated energy release are governed
primarily by the available reaction extent, reacting-material content,
reaction enthalpy, and, for the cathode, the relative contributions of
the two kinetic channels.

The cathode reaction deserves particular attention because its two
parallel kinetic channels share one conversion variable. Its integrated
energy therefore depends on the temperature-dependent partitioning of
reaction progress between the two channels and cannot generally be
represented by a single constant reaction enthalpy.


\subsection{Thermal response}

The combined reaction network releases

\begin{equation}
    q_{\mathrm{reaction}}
    =
    \SI{\HuIntegratedEnergy}{\kilo\joule\per\kilogram}
\end{equation}

during the simulated interval.

This heat release increases the lumped cell temperature from

\begin{equation}
    T_0
    =
    \SI{480}{\kelvin}
\end{equation}

to

\begin{equation}
    T_f
    =
    \SI{\HuFinalTemperature}{\kelvin}.
\end{equation}

The temperature evolution in
\cref{fig:hu_temperature_result}
reflects the positive feedback between Arrhenius reaction kinetics and
thermal energy release: increasing temperature accelerates the
exothermic reactions, which in turn generates additional heat.

The result should be interpreted as the response of the implemented
reduced-order model under the selected initial and boundary conditions,
rather than as a universal thermal runaway temperature for lithium-ion
cells.

The instantaneous thermal-power balance is shown in
\cref{fig:hu_thermal_power_balance}.

\begin{figure}[htbp]
    \centering
    \includegraphics[
        width=0.82\textwidth
    ]{hu2020_thermal_power_balance.pdf}

    \caption{
        Instantaneous mass-specific thermal-power balance for the
        Hu et al.\ reference simulation. Reaction heat generation is
        partitioned between sensible heating of the cell and convective
        heat transfer to the surroundings.
    }
    \label{fig:hu_thermal_power_balance}
\end{figure}

\subsection{Integrated reaction energy}

The total specific reaction energy is calculated from

\begin{equation}
    q_{\mathrm{reaction}}
    =
    \int
    \dot{q}_{\mathrm{total}}
    \,dt.
    \label{eq:results_integrated_heat}
\end{equation}

For the present reference calculation, the integrated reaction energy is
approximately

\begin{equation}
    q_{\mathrm{reaction}}
    =
    \SI{\HuIntegratedEnergy}{\kilo\joule\per\kilogram}.
\end{equation}

The final simulated cell temperature is approximately

\begin{equation}
    T_f
    =
    \SI{\HuFinalTemperature}{\kelvin}.
\end{equation}

These values characterize the selected numerical reference case and
should not be interpreted as general thermal runaway characteristics of
lithium-ion cells.